%!TEX root = ../hutbthesis_main.tex

\chapter{绪论}

\section{研究背景与研究意义}

固定翼飞行器着陆阶段是属于飞行全过程中里风险相对较高的环节之一。与跟巡航阶段相比较而言，着陆的时候飞行高度较低、，速度变化较快、，操纵时间较短，容极易受到侧风、低空湍流、地面效应以及、姿态扰动等因素影响的作用，进从而导致引发接地点偏差、下沉率过大、姿态超限和、滑跑偏移等问题。若要是控制不得当的话，不仅但会影响着陆质量，还可能对给飞行器结构安全和、后续运行造成带来不利影响。

随着固定翼飞行器特别是无人机，于巡检、测绘、侦察、运输等诸多领域被广泛运用，飞行任务所处的环境变得越来越复杂，这对着陆控制系统的安全性、稳定性、适应性，提出了更为严苛的要求。当前现有的着陆控制研究，主要聚焦于控制算法、轨迹跟踪、姿态调节等方面，虽说能够在一定程度上提高着陆精度，然而对于操纵者在着陆进程里的信息获取、风险判断、操作辅助，所给予的关注却相对欠缺。

在人机协同日益强化的背景状况下，把人机交互强化思想引入到固定翼飞行器着陆控制系统当中，借助状态提示、风险告警、操纵辅助、安全约束等多种方式，能够帮助操纵者更为精准地把握飞行状态，及时对操作予以修正，进而降低着陆风险。所以，对基于人机交互强化的固定翼飞行器着陆控制系统展开研究，具备颇为显著的现实意义。

从理论意义来讲，展开对于基于人机交互强化的固定翼飞行器着陆控制系统的研究，能够推动着陆控制研究从单一控制律设计朝着``控制层---交互层---保护层''协同设计方向延伸拓展，进而丰富固定翼飞行器着陆控制系统的研究内容，并且完善相关评价指标模式、设计方法。

从工程意义方面分析，此项研究可为固定翼飞行器在复杂工况时实现安全着陆提供更具操作性的系统方案，能为后续着陆控制仿真验证、模块化设计、工程应用给予参考。以人机交互强化作为切入点来研究固定翼飞行器着陆控制系统，具备较强的理论价值与实际应用价值。

\section{国内外研究现状}

外国对于固定翼飞行器着陆的研究重点放在起落架结构设计、低速着陆气动特性、缓冲系统动力学、控制辅助等领域。Elahi等人针对无人作战飞行器CFRP起落架展开设计与分析，觉得复合材料起落架能在保证承载能力的情况下达成轻量化，为着陆阶段结构安全设计给出了参照。

Lampropoulos等人针对低速无人机展开研究，聚焦于其在起飞、着陆、大气湍流状况下的机翼气动性能表现。他们发现低空区域存在的湍流、低速流场发生的变化，会对飞行器的升阻特性、姿态稳定性产生明显影响，进而对着陆控制精度造成影响。Dinc等人则对舰载无人机油气式起落架系统的动力学特性做了分析，探究出气体压缩、液体节流、结构非线性耦合作用对于缓冲响应与冲击载荷具有重要作用。Jinjin等人围绕无人机起落架缓冲器，从优化设计、试验验证这两方面展开研究，其结果显示对结构参数与阻尼特性进行合理优化，能够有效降低着陆冲击峰值，同时提高缓冲性能。

另外，Baek等人提出了基于输入信号变换的控制辅助系统，目的在于改善无人机目标跟踪与着陆的性能，这为把人机交互强化思想引入着陆控制系统给予了启发\linkcite{5}。从整体情况来看，国外的研究成果比较丰富，不过多数研究仍侧重于结构、缓冲或控制中的某一项内容，对于控制、交互与安全约束协同设计的系统性研究相对欠缺。

在国内，针对固定翼飞行器、无人机着陆展开的研究里，重点聚焦于起落架结构的优化工作、整体布局的分析工作、缓冲器性能的研究工作等多个方面。帅涛针对某一型号的起落架进行了设计、模态分析，明确指出起落架的固有频率、振型对于着陆冲击响应会产生非常重要的影响。周著学、于萍借助有限元分析方法针对小型航拍无人机的起落架进行了优化设计，阐述了合理调整结构、材料参数能够让应力分布得到改善并且提高承载能力。薛影、张超、吴星星剖析了无人机起落架布置对于起飞性能所造成的影响，显示出起落架的几何布局、安装位置与飞行性能存在着显著的耦合关系，而这对于着陆控制系统的总体设计同样具备参考价值。

罗杰等人针对某型无人机起落架放下阶段缓冲器的气液流动特性展开了研究，进而揭示出缓冲器内部气液流动变化对于阻尼性能、响应过程所产生的影响。蒲雅熙是从油气缓冲器气液两相阻尼特性着手，对不同工况下缓冲器阻尼变化规律予以分析，表明着陆速度与冲击条件会直接作用于缓冲器工作状态\linkcite{10}。

目前国内在起落架设计、有限元优化还有缓冲性能分析这些方面已经取得了一些成果。不过现有研究更多的是关注结构设计、缓冲机理，对于着陆阶段操纵者状态感知、风险提示、操作辅助与输入约束等人机交互强化问题关注不足。

整体看来，国内外研究已经给固定翼飞行器着陆控制奠定了不错的基础，特别是在起落架结构设计、低速气动特性、缓冲器动力学等方面取得了较多成果。然而现有研究大多是围绕结构、缓冲或者控制其中某一个方面进行的，针对``人机交互强化''的系统研究比较欠缺，特别是在复杂工况下怎样借助状态提示、风险告警、操纵辅助、输入约束来提高着陆安全性与稳定性，依旧缺少比较完整的研究框架。所以本文在已有研究的基础上引入人机交互强化思想，探究固定翼飞行器着陆控制系统的优化设计，具备一定的理论意义与实际价值。

\section{研究内容与技术路线}
\subsection{研究内容}
本文以固定翼飞行器着陆过程为研究对象，围绕人机交互强化的着陆控制系统开展研究。首先，对固定翼飞行器着陆任务场景、系统组成及关键影响因素进行分析，明确着陆过程中速度、高度、下沉率、姿态偏差等主要参数及其作用。其次，建立着陆阶段动力学模型和仿真模型，为后续分析不同工况下的着陆性能提供基础。再次，设计模型校核方法、性能评价指标和典型工况，通过仿真计算分析不同工况下系统响应变化规律，识别不利工况与较优工况。最后，结合仿真结果提出人机交互强化策略，重点从风险告警、状态提示、操纵辅助和约束控制等方面进行分析，并提出系统优化建议。
\subsection{技术路线}
本文围绕着``需求分析---模型建立---仿真验证---结果分析---策略优化''这一思路来进行研究工作。一开始，借助查阅相关文献，同时紧密结合课题任务要求，去剖析固定翼飞行器着陆阶段的主要风险、控制需求，进而确定研究内容和评价指标。接着，依据着陆过程的特点建立动力学模型与仿真模型，并完成主要参数的设定工作。然后，设计不同工况下的仿真方案，针对着陆过程里的速度、姿态、下沉率、偏差等指标展开计算与比较分析。最后，基于上述工作识别系统存在的问题，提出人机交互强化策略，并且对其工程可实现性予以讨论，从而形成论文的研究结论。
